% Source: Notion page “14. 项目完成情况”. Generated without rewriting its content.
\chapter{项目完成情况}

本项目围绕校园场景下的六类行为识别任务，构建了“骨架特征提取—时空关系建模—大模型语义推理—伪标签筛选—知识蒸馏—难例挖掘—增量学习—边缘部署”的完整技术链路。项目已完成全部必选功能，并进一步实现了Web可视化、人工审核、模型发布与回滚等工程化功能，形成可解释、可进化、轻量化且兼顾隐私保护的行为识别系统。

\section{功能模块完成情况}

\begin{center}
\small
\begin{longtable}{|>{\raggedright\arraybackslash}p{0.30\linewidth}|>{\raggedright\arraybackslash}p{0.30\linewidth}|>{\raggedright\arraybackslash}p{0.30\linewidth}|}
\hline
模块 & 项目实现方式 & 完成情况 \\
\hline
时空特征提取 & 原始视频首先经过RTMDet完成人体检测，再利用RTMPose提取人体关键点，统一转换为COCO-17骨架序列。进一步从连续骨架中计算人员中心位置、运动速度、轨迹变化、姿态覆盖率、双人距离等可量化时空特征，为后续行为识别和语义推理提供结构化输入。 & 已完成 \\
\hline
时空图构建 & 基于连续骨架轨迹构建Spatio-Temporal Graph。以场景中的人员作为节点，以人员之间的接近、距离变化、可能接触等关系作为边，并记录关系持续时间、最小距离、运动状态及姿态质量等信息，将原始坐标序列提升为可表达交互关系的时空结构。 & 已完成 \\
\hline
语义映射层 & 设计结构化语义映射模块，将时空图中的人员状态、轨迹特征、交互关系、时间片段和姿态质量整理为标准化Semantic Graph，并自动转换为版本化Prompt。Prompt中明确限定六类行为标签，并要求大模型仅依据实际测量证据进行判断，禁止臆测攻击意图、恐惧、防御等无法从骨架直接观测的信息。 & 已完成 \\
\hline
大模型推理模块 & 引入Qwen3-VL作为语义Teacher，对匿名骨架视频及结构化时空语义信息进行多模态推理。模型输出六类行为概率分布、置信度、判断依据、反证信息及是否需要人工审核等结构化结果，实现对传统骨架模型难例的语义补充判断。 & 已完成 \\
\hline
伪标签筛选策略 & 并非直接采用大模型输出作为训练标签，而是建立多信号质量评分机制，综合大模型多次推理一致性、Teacher置信度、Student-Teacher一致性、姿态质量及时序稳定性进行评分。同时针对师生高置信冲突、推搡缺少接触证据、追逐缺少持续相对运动、骨架跟踪异常等情况自动进入人工复核队列，保证伪标签质量。 & 已完成 \\
\hline
知识蒸馏框架 & 建立ProtoGCN Teacher-Student蒸馏框架。最终部署模型采用Logits Knowledge Distillation，通过真实标签交叉熵与Teacher软概率分布的KL散度联合训练，使轻量Student继承高精度模型的类别分布信息；代码中同时实现了可选择的Early/Middle/Deep三阶段特征对齐机制，可进一步进行多层级蒸馏。 & 已完成 \\
\hline
增量学习机制 & 建立完整的数据闭环。系统根据Student熵值、Top-1/Top-2间隔、时序不稳定性、师生分歧、低频类别等因素进行难例挖掘；难例经过Qwen辅助判断和人工审核后进入增量数据池。增量训练采用“新增难例+类别均衡历史样本”的Replay策略，并设置候选模型评测、发布门控和模型回滚机制，避免新增数据导致旧类别灾难性遗忘。 & 已完成 \\
\hline
可视化展示 & 实现基于FastAPI和Web前端的行为识别可视化系统，可展示六类行为预测结果、Top-K概率、Skeleton视频、姿态质量、难例队列、大模型推理依据、人工审核记录、模型状态及增量训练状态等。同时提供模型发布、审核和系统运行状态相关接口，实现从算法到系统的完整演示。 & 已完成（可选项已实现） \\
\hline
\end{longtable}
\end{center}

\section{综合技术要求完成情况}

\begin{center}
\small
\begin{longtable}{|>{\raggedright\arraybackslash}p{0.30\linewidth}|>{\raggedright\arraybackslash}p{0.30\linewidth}|>{\raggedright\arraybackslash}p{0.30\linewidth}|}
\hline
类别 & 项目完成情况 & 满足情况 \\
\hline
可解释性 & 系统不仅输出最终行为类别，还输出结构化推理依据。大模型结果包含行为概率分布、置信度、evidence、counter\_evidence、reason及needs\_review等字段；其依据直接来自轨迹速度、人员距离、接近过程、姿态覆盖率等可测量信息。例如，可通过“双人持续接近且距离显著缩短”“轨迹存在持续跟随关系”等证据解释追逐类判断，从而避免传统分类模型仅给出标签而无法说明原因的问题。 & 满足 \\
\hline
可进化性 & 构建“Student预测→难例发现→大模型辅助判断→人工审核→伪标签/真实标签沉淀→增量训练→候选模型评测→重新部署”的闭环机制。系统可持续收集真实场景中的低置信和高价值样本，并利用历史样本Replay维持原有能力，使模型能够随着新数据持续优化。系统还设计了每轮“50条新难例+250条类别均衡历史样本”的300条增量训练策略。 & 满足 \\
\hline
轻量化 & 在高精度FP32 ProtoGCN基础上采用QAT量化感知训练和Logits知识蒸馏得到M1KD INT8学生模型。最终模型文件约5.4 MiB，远低于题目要求的50 MB，同时在358条Campus6基线数据上取得345/358的整体正确结果，实现模型精度与边缘部署体积之间的平衡。 & 显著满足 \\
\hline
鲁棒性 & 输入侧采用人体骨架而非RGB外观作为主要识别信息，从机制上降低服装颜色、背景和光照变化对行为分类的干扰；训练阶段采用时间均匀采样、左右翻转等数据增强，并统一转换至COCO-17骨架空间。系统同时监测姿态覆盖率、跟踪失败和时序稳定性，对遮挡严重或姿态质量较低的样本不进行盲目高置信判断，而是自动进入人工复核流程，从而提升复杂场景下的可靠性。 & 满足 \\
\hline
隐私保护 & 系统采用Pose-only隐私策略。原始视频仅在边缘侧短暂用于人体检测与姿态估计，随后转换为匿名COCO-17骨架数据；Web端仅展示Skeleton视频及结构化语义信息，不向浏览器传输原始RGB视频。配置中关闭RGB回退，并设置原始视频保留时间为0，从数据链路上减少人脸、服装、身份及场景隐私泄露风险。 & 满足 \\
\hline
代码规范 & 项目采用标准Python src工程结构，将后端服务、模型推理、训练流水线、语义Teacher、特征提取、评估、配置、可视化和测试进行模块化拆分；同时提供README、依赖文件、启动脚本、模型MANIFEST、设计及验收文档，并配套单元测试和集成测试。训练、推理、评估、难例分析等功能均提供独立命令行入口，便于复现、部署和后续维护。 & 满足 \\
\hline
\end{longtable}
\end{center}

\section{项目整体完成度总结}

综上，本项目已完成题目要求的全部核心模块，并在基本功能之外实现了较完整的工程化闭环。系统以前端RTMDet/RTMPose实现隐私友好的骨架感知，以ProtoGCN作为高效的六分类Student模型，以Qwen3-VL作为仅在难例条件下触发的语义Teacher，通过时空图和结构化Prompt提升模型的可解释能力；再结合多信号伪标签筛选、人工复核、知识蒸馏、难例挖掘和增量学习，使模型具备持续进化能力。最终通过QAT INT8与Logits蒸馏将部署模型压缩至约5.4 MiB，并配套Web可视化、审核追踪、模型发布和回滚机制，实现了从算法设计、模型训练到实际部署和持续迭代的完整校园行为识别方案。

